\section{Proofs for selections, crowds, entropy and deletion}\label{finapp:all}
The main text keeps the short prime-resource arguments beside their
conclusions.  This appendix supplies the quantitative estimates and
limit transitions used there.

\subsection{The binomial approximation}\label{finapp:binomial}
\begin{proof}[Proof of Corollary~\ref{fin:binomialtv}]
Write $M_m=\max_j\Prob(\operatorname{Binomial}(m,\theta)=j)$.
Unimodality gives
$d_{\mathrm{TV}}(Z+a,Z+b)\le |a-b|M_m$ for integer shifts.
Couple $R$ with an independent $\operatorname{Binomial}(d_B,\theta)$
and retain the common independent $Z$.  The distance is at most
$d_B M_{m_B}=O_\theta(d_B/\sqrt{m_B})$.
The prime number theorem gives
$r_B\le\pi(\sqrt X)=O(\sqrt X/\log X)$ and
$N_B\sim X/\log X$, proving the bound.  The other assertions follow
from the polynomial factorization.
\end{proof}

\subsection{The compulsory part of a random selection}\label{finapp:greedy}
We use Buchstab's function $\omega$: it satisfies $u\omega(u)=1$ for
$1\le u\le2$ and $(u\omega(u))'=\omega(u-1)$ for $u>2$.
The classical rough-number estimate is
\begin{equation}
 \Phi(x,y):=\#\{n\le x:\lpf(n)>y\}
 \sim \omega(u)\frac{x}{\log y},\qquad
 u=\frac{\log x}{\log y}>1
 \label{fin:rough}
\end{equation}
when $u$ stays near a fixed number greater than one; $1$ is counted.
We also use $\Phi(x,y)\ll x/\log y$ for $x\ge y\ge2$.
These forms follow, for example, from the estimates reviewed and made
explicit in \cite{finFanRough}.

\begin{proof}[Proof of Theorem~\ref{fin:greedy}]
Put $h=\log N$, $X=\lfloor N^\beta\rfloor$, and let $R$ count sampled
integers with no prime divisor at most $N$.  By \eqref{fin:rough} and
binomial concentration,
$R\sim\omega(\beta)N/h$ in probability.

A rough label fails to belong to $\mathcal U_N$ only by sharing a prime
$p>N$ with another label.  If $r_p$ is the probability that the first
label is rough and divisible by $p$, the expected number of such failures
is at most $N^2\sum_{p>N}r_p/p$.
For $p\le X/N$, the rough-number bound gives
$r_p\ll1/(p\log N)$; for $p>X/N$, the only rough cofactor is $1$,
so $r_p=1/X$.  Thus the expected number is
\[
 O_\beta\left(\frac{N}{h^2}+\frac{N^2}{X}\right)=o(N/h).
\]
Here $\sum_{p>N}p^{-2}\ll1/(N\log N)$ and
$\sum_{N<p\le X}p^{-1}=O_\beta(1)$.
It follows that $|\mathcal U_N|\ge R-o_{\Prob}(N/h)$.

For the matching upper bound, put $y=N/h^2$.  In any coprime selection,
each non-$y$-rough member uses a distinct prime at most $y$.
Its size is therefore at most $R_y+\pi(y)$, where $R_y$ is the number
of $y$-rough sampled labels.  Again \eqref{fin:rough} and binomial
concentration give $R_y\sim\omega(\beta)N/h$, while
$\pi(y)=o(N/h)$.  Since $|\mathcal U_N|\le\alpha_N$, both asymptotics
follow.

Finally every maximal selection contains $\mathcal U_N$, so
\[
 |S\mathbin{\triangle}T|\le2(\alpha_N-|\mathcal U_N|).
\]
A greedy scan always produces a maximal selection.  All conclusions
hold simultaneously over its possible orders.
\end{proof}

\subsection{Crowds, missing primes and repeated powers}\label{finapp:crowds}
\begin{proof}[Proof of Proposition~\ref{fin:coverage}]

In the conditioned experiment,
\[
 \Prob(L_n>y)=\prod_{p\le y}\frac n{n+p-1}.
\]
For $y=t\sqrt{n\log n}$, the negative logarithm equals
$n^{-1}\sum_{p\le y}(p-1)+
O(n^{-2}\sum_{p\le y}p^2)\to t^2$ by the prime number theorem.

In the ordinary experiment the survival probability is
$\prod_{p\le y}(1-(1-1/p)^n)$.
Put $a=\log n-3\log\log n+t$ and $y=n/a$.
The quantitative prime number theorem and partial summation give
\[
 \sum_{p\le y}(1-1/p)^n
 \sim \int_2^y\frac{e^{-n/u}}{\log u}\,du
 =n\int_a^{n/2}\frac{e^{-v}}{v^2\log(n/v)}\,dv
 \sim\frac{ne^{-a}}{a^2\log(n/a)}\to e^{-t}.
\]
Primes below $y/2$ contribute $o(1)$; on $[y/2,y]$ replacement by the
integral is valid even on its effective width $y/\log n$, using the
standard prime-number-theorem error $O(y e^{-c\sqrt{\log y}})$.
The largest summand tends to zero, so the product tends to
$e^{-e^{-t}}$.  This is the asserted law after the change $y=n/a$.
\end{proof}

The following local calculation also proves the remainder law
\eqref{fin:remainderlaw}.
To verify it, the excess exponent at $p$ has the representation
$B_pG_p$, where $B_p$ is Bernoulli with parameter $n/(n+p-1)$ and
$G_p$ is geometric with $\Prob(G_p=j)=(1-1/p)p^{-j}$, independently.
The nonzero excesses occur at only finitely many primes almost surely.
Taking the ratio of the mass at $j$ to the mass at zero gives
$(n/(n+p))p^{-j}$ and proves the formula.

\begin{proof}[Proof of Corollary~\ref{fin:dickman}]
Couple the preceding excess exponents with
$H_n=\prod_{p\le n}p^{G_p}$.  The chance that any exponent differs is
at most
$\pi(n)/n+n\sum_{p>n}p^{-2}=O(1/\log n)$.
For $H_n$, logarithmic Laplace transforms expand into
\[
 \sum_{p\le n}\sum_{j\ge1}\frac{p^{-j}}j
       (e^{-sj\log p/\log n}-1).
\]
The terms $j\ge2$ vanish by dominated convergence.  Mertens' theorem
turns the $j=1$ terms into the asserted integral.

The limit is the sum of a Poisson process on $(0,1]$ with intensity
$dt/t$.  Its largest point is uniform, so
$Y\overset d=U(1+Y')$, with $U$ uniform and $Y'$ an independent copy.
Its density equals $\E[(1+Y')^{-1}{\bf1}_{y<1+Y'}]$, constant for
$0<y<1$.  The Laplace transform is asymptotic to $e^{-\gamma}/s$ as
$s\to\infty$, identifying that constant as $e^{-\gamma}$.
\end{proof}

\subsection{Real roots and the averaged Wishart polynomial}\label{finapp:matrix}
For completeness, each prime acts on a polynomial by positive rescaling
and $P\mapsto P+a zP'$, $a>0$.  This preserves negative real roots:
for $P(z)=\prod_j(z+r_j)$, $r_j>0$, the imaginary part of
$zP'(z)/P(z)=\sum_j z/(z+r_j)$ is strictly positive in the upper
half-plane.  The new polynomial has no roots there or in the lower
half-plane, and its positive coefficients exclude nonnegative roots.
Start with $(1+z)^N$ and pass to the coefficient limit over primes.

\begin{proof}[Proof of Theorem~\ref{fin:wishart}]
Cauchy--Binet expresses a principal minor of $Z_NZ_N^*$ of order $j$
as a sum of $\binom Nj$ squared Gaussian determinants, each with mean
$j!$.  Consequently, for the Laguerre polynomial
$L_N(x)=\sum_k\binom Nk(-x)^k/k!$,
\[
 \E\det(Z_NZ_N^*+tI)=N!L_N(-t).
\]
Since $\E \Comp^k=k!K(k)$, averaging at $t=\Comp z$ gives $N!J_N(z)$.
Power biasing by $\Comp^N$ proves the characteristic-polynomial identity.

Put $\kappa(t)=-\log K(t)$.  Differentiating its Euler product and
using the prime number theorem and Mertens' product theorem gives
\[
 \kappa'(t)=\log\log t+\gamma+o(1),\qquad
 \kappa''(t)\sim\frac1{t\log t}.
\]
For the second estimate, substitute $u=tv$ in
$\int_2^\infty(u+t-1)^{-2}\,d\pi(u)$; for the first, split the
logarithmic derivative at $p=t$.  The correction to
$-\sum_{p\le t}\log(1-1/p)$ is $O(1/\log t)$ by partial summation.
Write $b_N=K(N-1)/K(N)\sim a_N$.  Coefficient comparison yields
\[
 \frac1N\sum_jR_{N,j}=b_N,\qquad
 \frac1N\sum_j(R_{N,j}-b_N)^2
 =(N-1)\left(b_N^2-\frac{K(N-2)}{K(N)}\right)
 \sim\frac{b_N^2}{\log N}.
\]
The last equivalence uses the second difference of $\kappa$, asymptotic
to $1/(N\log N)$.  Chebyshev's inequality gives the root limit.

The biased scalar satisfies
\[
 \E \CompN =(N+1)\frac{K(N+1)}{K(N)}
       \sim e^{-\gamma}N/\log N,\qquad
 \frac{\Var(\CompN )}{(\E \CompN )^2}
 =\frac{N+2}{N+1}\frac{K(N+2)K(N)}{K(N+1)^2}-1\sim\frac1N.
\]
It therefore concentrates at that scale.  Apply the square
Marchenko--Pastur theorem \cite{finMarchenkoPastur} to
$Z_NZ_N^*/N$ and multiply by $N/(\CompN a_N)\to1$.
\end{proof}

\subsection{The prime-factor coupling and its moments}\label{finapp:entropy-coupling}
Write $P=(P_i)$ for the limiting partition, and set
$H=-\sum_iP_i\log P_i$, $D=e^H$.  The correlation formula
\cite[Theorem 2.1]{finHanda} is
\begin{equation}\label{fin:pdcorrelation}
 \E\sum_{i_1,\ldots,i_k\ {\rm distinct}}
 F(P_{i_1},\ldots,P_{i_k})
 =
 \int_{\substack{u_j>0\\\sum u_j<1}}
 F(u_1,\ldots,u_k)\prod_j\frac{du_j}{u_j}.
\end{equation}
One can see this formula directly from the gamma process: take Poisson
jumps with intensity $e^{-v}dv/v$, divide by their exponential total,
and use independence of that total and the normalized jumps.
After choosing $k$ jumps and writing the leftover total as
$s(1-\sum u_j)$, the Jacobian cancels their common power of $s$;
integrating $e^{-s}ds$ leaves \eqref{fin:pdcorrelation}.
Taking $k=1$ and $F(u)=u\log(1/u)$ also gives
$\E H=\int_0^1\log(1/u)\,du=1$.

\begin{lemma}[Grouping repeated primes]\label{fin:entropycoupling}
For a uniform integer $U_N\in[1,N]$, let $Q_N$ be its decreasing
prime-power share sequence, with $Q_N=(1,0,\ldots)$ at $U_N=1$.
There is a coupling with $P$ such that
\[
 \E\|Q_N-P\|_1\ll(\log N)^{-1},\qquad
 \E|H(Q_N)-H(P)|\ll_\beta(\log N)^{-\beta}
 \quad(0<\beta<1).
\]
Moreover $\sup_N\E e^{tH(Q_N)}<\infty$ for every fixed $t>0$.
\end{lemma}
\begin{proof}
Haddad and Koukoulopoulos \cite{finHaddadKoukoulopoulos} couple the
decreasing prime factors $q_i$ \emph{with multiplicity} to $P$ so that
$\E\sum_i|\log q_i-P_i\log N|=O(1)$.
Changing the denominator from $\log N$ to $\log U_N$ costs
$1-\log U_N/\log N$ in $\ell^1$, with mean $O(1/\log N)$.
The exceptional value $U_N=1$ costs $O(1/N)$.

Grouping repeated occurrences of $p$ changes the vector by at most
twice their excess logarithmic mass
\[
 R(n)=\frac1{\log n}\sum_p\sum_{j\ge2}
                  (\log p){\bf1}_{p^j\mid n}.
\]
For $n\ge\sqrt N$, divisibility counting bounds its mean by
$2(\log N)^{-1}\sum_p(\log p)/(p(p-1))$.
The remaining integers contribute at most $N^{-1/2}$.
Sorting minimizes matching distance and cannot increase this bound.
This proves the first assertion.

We record a moment bound used twice.  For $0<\alpha<1$ and a mass
partition $Q$, write $S_\alpha(Q)=\sum_iQ_i^\alpha$.
For every fixed positive integer $r$,
\begin{equation}\label{fin:powerbound}
 \sup_N\E S_\alpha(Q_N)^r<\infty.
\end{equation}
On $n\ge\sqrt N$, bound $S_\alpha$ by
$\sum_{p^j\mid n}(2\log p/\log N)^\alpha$ and expand its $r$th power.
A block of $b$ equal primes has exponent sum
\[
 \sum_{j_1,\ldots,j_b\ge1}p^{-\max(j_1,\ldots,j_b)}
 =\sum_{m\ge1}\frac{m^b-(m-1)^b}{p^m}\ll_b\frac1p.
\]
Then $\sum_{p\le N}(\log p)^{\alpha b}/p
\ll_{\alpha,b}(\log N)^{\alpha b}$ bounds each block, and there are
only finitely many partitions of $r$ occurrences.
Integers below $\sqrt N$ contribute negligibly because
$S_\alpha\le(\log N/\log2)^{1-\alpha}$.

For $h(u)=u\log(1/u)$, $\alpha=(1-\beta)/2$, elementary differentiation
and a split into $u\le v/2$ and $v/2<u\le v$ give
\[
 |h(u)-h(v)|\ll_\beta |u-v|^\beta(u^\alpha+v^\alpha).
\]
Hölder's inequality, first on coordinates and then on expectations,
bounds the entropy difference by a constant times
$(\E\|Q_N-P\|_1)^\beta$; the remaining factors are bounded by
$\E S_{1/2}(Q_N)$ and
$\E S_{1/2}(P)=2$, the latter from \eqref{fin:pdcorrelation}.
Finally Jensen's inequality gives
$e^{(1-\alpha)H(Q)}\le S_\alpha(Q)$, so
\eqref{fin:powerbound} bounds every fixed exponential moment.
\end{proof}
It follows that the limiting mean effective count is
$\mathfrak P=\E e^H$, numerically about $2.9621$.
This is different from $e^{\E H}=e$.
The following recursion makes that numerical constant reproducible.
If $h_2(u)=-u\log u-(1-u)\log(1-u)$, stick breaking gives
$H\overset d=h_2(U)+(1-U)H'$.
For $a_r=\E H^r/r!$,
\[
 a_0=1,\qquad
 a_r=\frac{r+1}{r}\sum_{j<r}a_j
       \int_0^1\frac{h_2(u)^{r-j}}{(r-j)!}(1-u)^j\,du,\qquad
 \mathfrak P=\sum_{r\ge0}a_r.
\]
All terms are nonnegative.  Writing the entropy as a sum of
stick-breaking contributions bounds
$H\le(\log2)(1+\sum_{j\ge1}\prod_{\ell\le j}(1-U_\ell))$.
The product terms are exponentials of a rate-one Poisson process's
arrival times; its exponential formula gives
$\E e^{tH}\le2^t\exp(2^t-1)$.
Thus truncating the last series at $r=R$ has error at most
$2^{-R-1}\E e^{2H}<81\,2^{-R-1}$.
Lemma~\ref{fin:entropycoupling}, with uniform integrability, also proves
$\E D(U_N)\to\mathfrak P$.


\subsection{The sharp entropy tail}\label{finapp:entropy-tail}
\begin{proof}[Proof of Theorem~\ref{fin:entropytail}]

Let $M(t)=\E e^{tH}$ and $h(u)=u\log(1/u)$.
Expanding a nonnegative product and using
\eqref{fin:pdcorrelation} gives, for every $a>0$,
\[
 \log M(t)\le a+I_t(a),\qquad
 I_t(a)=\int_0^1 e^{-au}(e^{th(u)}-1)\frac{du}{u}.
\]
Here the simplex constraint was bounded by
$e^{a(1-\sum u_j)}$.  Put $L=\log t$, choose $b>0$ from
$b+\tfrac12\log b+\tfrac12\log(2\pi)=L$, and set
$a=t(\log(t/b)-1)$.
Then
\[
 b=L-\tfrac12\log L-\tfrac12\log(2\pi)+o(1).
\]
Under $z=tu/b$, the positive exponent becomes
$b\chi(z)$ with $\chi(z)=z(1-\log z)$, whose unique maximum is $\chi(1)=1$
and whose second derivative there is $-1$.
Laplace's method yields
\[
 I_t(a)=e^b\sqrt{2\pi/b}(1+o(1))
       =\frac tb(1+o(1)).
\]
To control the endpoint, the cancellation in $e^{th(u)}-1$ makes the
integral over $0<z<b^{-2}$ bounded.  On
$[b^{-2},1/2]\cup[2,e]$ the maximum of $\chi$ is strictly below one,
and on $[e,t/b]$ the integral is $O(\log t)$.
All these errors are $o(e^b/\sqrt b)$.
Consequently
\[
 \log M(t)\le t(L-\log L-1)
 +\frac tL\left(\tfrac12\log L+1+\tfrac12\log(2\pi)\right)
 +o(t/L).
\]
Markov's inequality with $t=x\log x$ gives the required upper bound
on the probability.

For the reverse bound write $X=\log x$ and choose
\[
 \eta=X^{-1/8},\quad \epsilon=X^{-1/4},\quad
 v=(1-\eta)/X,\quad
 k=\left\lceil x\exp((1/2+\eta)/X)\right\rceil .
\]
In $\mathbb R^k$ consider the region
\[
 \mathcal R=\left\{ |z_i|\le\epsilon,\quad
  -\tfrac12\le\sum_i z_i\le0,\quad
  (1-\eta/4)kv\le\sum_i z_i^2\le(1+\eta/4)kv\right\}.
\]
Set $u_i=(1+z_i)/k$.  Their total lies in $[1-1/(2k),1]$.
The unlisted mass is too small to contain another component in
$[(1-\epsilon)/k,(1+\epsilon)/k]$.
Thus the factorial count, divided by $k!$, is an event probability,
rather than an overcount.

Taylor's inequality
$(1+z)\log(1+z)\le z+(\tfrac12+C\epsilon)z^2$ shows that on this event
\[
 H\ge(1-\tfrac1{2k})\log k
       -(\tfrac12+C\epsilon)(1+\eta/4)v>\log x
\]
for large $x$.  Also
$\prod_i(1+z_i)^{-1}\ge1$, so
\[
 \Prob(D>x)\ge\frac{\operatorname{vol}(\mathcal R)}{k!}.
\]

We give the volume estimate to identify the constant.
Take independent $N(0,v)$ variables, each conditioned to lie in
$[-\epsilon,\epsilon]$.  The conditioning costs
$\exp(-o(k))$, changes the variance by a relative $o(\eta)$, and
keeps the standardized third moments bounded.
Berry--Esseen gives probability at least $c/\sqrt{kv}$ for their
sum to lie in $[-1/2,0]$: its $O(k^{-1/2})$ error is smaller than this
slab probability.  Chebyshev bounds failure of the quadratic shell by
$O((k\eta^2)^{-1})=o((kv)^{-1/2})$.
Undoing the Gaussian density on the shell therefore gives
\[
 \log\operatorname{vol}(\mathcal R)
 \ge \tfrac k2\log(2\pi ev)-o(k).
\]
Stirling's formula now yields
\[
 -\log\Prob(D>x)
 \le k\log k-k-\tfrac k2\log(2\pi ev)+o(k)
 =x\left(X+\tfrac12\log X-1-\tfrac12\log(2\pi)+o(1)\right).
\]
This matches the first bound.
\end{proof}

\subsection{The Gaussian profile and its transfer to integers}\label{finapp:entropy-profile}
\begin{proof}[Proof of Theorem~\ref{fin:entropyprofile}]

Put $X=\log x$, $t=xX$, and use $a,b$ from the preceding saddle
calculation.  Fix a continuous compactly supported function $\psi$,
and write $T_x=\int\psi\,d\nu_x$.
For fixed real $s$, the exponent
\[
 q(u)=tu\log(1/u)+sxu\,\psi(\sqrt X(xu-1))
\]
is nonnegative for large $x$.  On the perturbation's support,
$u=(1+O(X^{-1/2}))/x$, so its first term is of order $X^2$
and dominates the bounded perturbation.  The same nonnegative
factorial expansion gives
\[
 \E e^{tH+sxT_x}\le \exp(a+I_t(a)+\Delta_x(s)),
\]
where
\[
 \Delta_x(s)=\int_0^1 e^{tu\log(1/u)-au}
        \left(e^{sxu\psi(\sqrt X(xu-1))}-1\right)\frac{du}{u}.
\]
Under $y=\sqrt X(xu-1)$, uniformly on that fixed compact support,
the exponential is $e^{b-y^2/2+o(1)}$, and
$du/u=(1+o(1))dy/\sqrt X$.
The saddle equation gives $e^b/(x\sqrt X)\to1/\sqrt{2\pi}$.
Hence, with $Z$ standard Gaussian,
\[
 \frac{\Delta_x(s)}x\longrightarrow
 \Lambda_\psi(s):=\E(e^{s\psi(Z)}-1).
\]
The sharp lower and upper bounds just proved also give
$\log\Prob(H>X)=a+I_t(a)-tX+o(x)$.
For $s>0$, exponential Markov therefore implies
\[
 \limsup\frac1x\log
 \Prob(T_x\ge \E\psi(Z)+\delta\mid H>X)
 \le\Lambda_\psi(s)-s(\E\psi(Z)+\delta)<0
\]
when $s$ is small enough.  Negative $s$ gives the lower deviation.
Thus compactly supported continuous tests converge, with exponentially
small conditional deviation probabilities.
All measures have mass one and the limit has mass one, so this vague
convergence is weak convergence.  On each fixed standardized interval,
$xP_i=1+O(X^{-1/2})$, allowing weights $P_i$ to be replaced by $1/x$.
This proves the count assertion.
\end{proof}
\begin{proof}[Proof of Corollary~\ref{fin:integerentropy}]

Let $L=\log N$, choose $\beta\in(1-\epsilon/2,1)$, and first take
$\eta=(\log x)^{-2}$.  Lemma~\ref{fin:entropycoupling} bounds the
probability of entropy error exceeding $\eta$ by
$O_\beta(\eta^{-1}L^{-\beta})$.
Both comparison tails at $xe^\eta$ and $xe^{-\eta}$ are at least
$L^{-(1-\epsilon/2)+o(1)}$.
The coupling error is negligible relative to them.  Sandwiching the
integer event between these two tails proves the tail formula because
$\eta\log x\to0$ changes the exponent by only $o(x)$.

For the profile, choose instead
$\eta=x^{-2}(\log x)^{-2}$.  For a compactly supported Lipschitz $\psi$,
the function
$u\mapsto u\psi(\sqrt{\log x}(xu-1))$ is Lipschitz with constant
$O_\psi(1+x\sqrt{\log x})$.  Thus its coupled sums differ in expectation
by $O_\psi((1+x\sqrt{\log x})/L)$.
Both this Markov error and the entropy-error probability remain
negligible relative to the integer conditioning probability.
The bad-profile probability in the limiting model, on
$H>\log x-\eta$, is at most
$\exp(-c_\psi x+o(x))$ times its tail probability, by the proof of
Theorem~\ref{fin:entropyprofile}.  The shift changes the saddle
normalization by $o(x)$ and the standardized test by $o(1)$.
The ratio of that tail to the integer tail is $\exp(o(x))$, by the
already proved logarithmic asymptotics.  Therefore its conditional
bad-profile probability still tends to zero.  A countable determining
set of Lipschitz tests completes the proof.
\end{proof}

\subsection{Why every finite deletion depth is abundant}\label{finapp:deletion}
The iteration needs a lower bound on the \emph{number} of Goldbach
representations; mere almost-everywhere existence would not suffice.
\begin{lemma}[Robust almost-all Goldbach]\label{fin:robust}
Let $R_2(m)$ count ordered prime pairs summing to the even integer $m$.
There is $c>0$ such that, for every $A>0$,
\[
 \#\{m\le Y:m\text{ even},\
       R_2(m)<cm/(\log m)^2\}\ll_A Y/(\log Y)^A.
\]
\end{lemma}
\begin{proof}
We use the weighted form explicitly obtained in the proof of
\cite[Theorem 4.3]{finBhowmikGrimmelt}, the classical almost-all
Goldbach argument.  Outside a set with arbitrarily large logarithmic
saving, it gives
\[
 \sum_{a+b=m}\Lambda(a)\Lambda(b)
             \ge \tfrac12m\mathfrak S(m).
\]
The singular series has a positive absolute lower bound for even $m$.
Pairs involving a non-prime prime power contribute
$O(\sqrt m(\log m)^3)=o(m)$: there are
$O(\sqrt m\log m)$ possible proper prime powers and each weight product
is at most $(\log m)^2$.  Remove them and bound each remaining
prime-pair weight by $(\log m)^2$.
This gives the asserted lower bound, with a smaller $c$.
One may apply the weighted estimate on dyadic intervals and sum them;
the finitely many small targets are absorbed in the implicit constant.
\end{proof}
\begin{proof}[Finite-depth assertions of Theorem~\ref{fin:delete}]
For abundance, write $D_k(t)=\#\{p\le t:p\notin B_k\}$.
The estimate at $k=0$ is immediate when starting from all primes.
For the Goldbach-enabled initial set it follows from
Lemma~\ref{fin:robust}, since an index $i$ with $p_i\le X$ has target
at most $2(\pi(X)+1)$.

Assume the estimate at depth $k$.  Each missing prime spoils at most
two ordered pairs for a fixed target, so
\[
 R_{2,B_k}(m)\ge R_2(m)-2D_k(m).
\]
Taking the induction estimate with exponent $4$ gives
$D_k(m)=o_k(m/\log^2m)$.
Every sufficiently large robust target from Lemma~\ref{fin:robust}
therefore still has a certificate at depth $k$.
The newly missing indices through $\pi(X)$ form a subset of the
lemma's exceptional targets through $2(\pi(X)+1)$.
Together with the previously missing primes and finitely many small
targets this proves the estimate at depth $k+1$, for every requested
logarithmic power.  Its constants can depend on $k$.

For the ternary statement, the same robust lemma already implies
$R_3(n)\gg n^2/\log^3n$ for odd $n$: take primes
$q\in[n/3,n/2]$.  The prime number theorem supplies
$\gg n/\log n$ choices, and all but $O_A(n/\log^A n)$ have robust
even target $n-q$.  Each such target has
$\gg n/\log^2n$ representations.
A triple using a missing prime can be counted by choosing its position,
the missing prime, and one other prime.  There are at most
$3D_k(n)\pi(n)=o(R_3(n))$ such triples.  Hence
\mbox{$R_{3,B_k}(n)/R_3(n)\to1$}.
\end{proof}

\subsection{The graph spectrum and its collapse}\label{finapp:spectrum}
The two-dimensional local matrix and its eigenvalues are
\[
 L_p=\begin{pmatrix}
 1-1/p&\sqrt{(1/p)(1-1/p)}\\
 \sqrt{(1/p)(1-1/p)}&0
 \end{pmatrix},\qquad
 a_p^\pm=\frac{p-1\pm\sqrt{(p-1)(p+3)}}{2p}.
\]
Put $b_p=-a_p^-/a_p^+$ and $\Lambda_B=\prod_{p\in B}a_p^+$.
The product is positive since $a_p^+=1+O(p^{-2})$.
The nonzero spectrum, with multiplicity, is
\begin{equation}\label{fin:localspectrum}
 \Lambda_B(-1)^{|F|}\prod_{p\in F}b_p,\qquad F\subset B\text{ finite}.
\end{equation}
For the full prime set this is Zagier's spectrum
\cite[Appendix C]{finCoprimeChain}.
The tensor calculation also proves the version for $B$.
Indeed finite prime cutoffs have exactly these products.
Their kernels converge in $L^2$ because
$\sum_{p>z}1/p^2\to0$; a product involving a moving prime tends to zero
since $b_p\sim1/p$.  Thus every eigenvalue bounded away from zero
comes from a fixed finite set $F$, establishing the entire nonzero
list.

\begin{lemma}[The spectral tail]\label{fin:weyl}
Let $B$ omit a prime set $D$ with $\sum_{p\in D}1/p<\infty$.
If $N_{B,+}(u)$ and $N_{B,-}(u)$ count eigenvalues of $T_B$ of each
sign and magnitude at least $u$, then
\[
 N_{B,+}(u)\sim N_{B,-}(u)\sim\frac{A_B}{2u},
\]
where
\[
 A_B=\sqrt{K_B(4)}\prod_{p\in D}(1-1/p),\qquad
 K_B(4)=\prod_{p\in B}(1-1/p)^4(1+4/(p-1)).
\]
\end{lemma}
\begin{proof}
For $B=\mathcal P$, the absolute and signed spectral Dirichlet series are
\[
 D_{\rm abs}(s)=\Lambda_{\mathcal P}^s\prod_p(1+b_p^s)
              =\zeta(s)\mathcal A(s),\qquad
 D_{\rm sgn}(s)=\Lambda_{\mathcal P}^s\prod_p(1-b_p^s)
              =\mathcal E(s)/\zeta(s).
\]
Since $b_p=p^{-1}+O(p^{-2})$, the correction $\mathcal A$ is analytic
for $\Re s>1/2$, and $\mathcal E$ is analytic for $\Re s>0$.
These assertions follow from locally uniform product estimates
$1+O(p^{-\Re s-1}+p^{-2\Re s})$ and
$1+O(p^{-\Re s-1})$, respectively.
At $s=1$,
\[
 \mathcal A(1)
 =\prod_p(1-1/p)(a_p^+-a_p^-)=\sqrt{K(4)}=:A.
\]
The identity follows by squaring each local factor.
The two positive-coefficient series
$(D_{\rm abs}\pm D_{\rm sgn})/2$ have a simple pole of residue $A/2$
at one and a continuous remainder on $\Re s=1$; the prime number
theorem supplies nonvanishing of $\zeta$ on that line.
The Wiener--Ikehara theorem \cite{finIkehara}, in its
Laplace--Stieltjes form, gives
$N_{\mathcal P,\pm}(u)\sim A/(2u)$.

For the omitted set $D$, divide out its local factors:
\[
 D_{{\rm abs},B}(s)
 =D_{\rm abs}(s)\Lambda_D^{-s}
                \prod_{p\in D}(1+b_p^s)^{-1},
 \quad
 D_{{\rm sgn},B}(s)
 =D_{\rm sgn}(s)\Lambda_D^{-s}
                \prod_{p\in D}(1-b_p^s)^{-1}.
\]
Expand each inverse factor into a geometric series.  At exponent one
the total absolute coefficient weight is bounded by
$\Lambda_D^{-1}\prod_{p\in D}(1-b_p)^{-1}<\infty$.
The full-prime counting functions are $O(x)$ at reciprocal threshold
$x$, so dominated convergence applies to these convolutions.
The absolute count has coefficient
\[
 \frac{A}{\Lambda_D\prod_{p\in D}(1+b_p)}
 =A\prod_{p\in D}\frac{p}{\sqrt{(p-1)(p+3)}},
\]
while the signed count remains $o(x)$.  The displayed coefficient
equals the asserted $A_B$, again by the local identity for $K(4)$.
\end{proof}
\begin{proof}[Proof of Theorem~\ref{fin:collapse}]
Theorem~\ref{fin:delete} and partial summation imply
$\sum_{p\notin B_k}1/p<\infty$ for each fixed $k$, so the spectral
tail lemma applies.  Its equivalent product is
\[
 A_k=\sqrt{K(4)}
       \prod_{p\notin B_k}\frac{p}{\sqrt{(p-1)(p+3)}}.
\]
Every factor is below one.  Their logarithms are
$-1/p+O(p^{-2})$, and the omitted sets increase to
$\mathcal P\setminus\{2,3,5\}$.  Divergence of the prime harmonic
series proves $A_k\downarrow0$.

Two independent prime profiles share only finitely many primes almost
surely, since $\sum_p1/p^2<\infty$.  All shared primes above $5$
eventually disappear from $B_k$.  Thus the kernels converge almost
everywhere to $G_{\{2,3,5\}}$, and bounded convergence gives
Hilbert--Schmidt convergence.
Its three local matrices each have two nonzero eigenvalues, so the
tensor has rank $2^3=8$.
Finally the local four-clique factors are below one, so removing
restrictions increases the product.  Its final value is
$(5/16)(16/27)(512/625)=512/3375$.
\end{proof}
